% =====================================================================
% AEGIS — artigo no modelo IEEE (conferência, duas colunas)
% Trabalho de Linguagens Formais e Compiladores
%
% Como usar no Overleaf: New Project > Upload Project > este arquivo
% (ou cole no main.tex de um projeto "IEEE Conference Template").
%
% Marcadores:
%   \pendente{...}  -> trecho que ainda falta escrever (sai em vermelho)
%   \figura{...}    -> lugar reservado para uma imagem sua; troque por
%                      \includegraphics quando subir o arquivo
% =====================================================================
\documentclass[conference]{IEEEtran}

\usepackage[brazil]{babel}
\usepackage[utf8]{inputenc}
\usepackage[T1]{fontenc}
\usepackage{cite}
\usepackage{amsmath,amssymb}
\usepackage{graphicx}
\usepackage{xcolor}
\usepackage{booktabs}
\usepackage{listings}
\usepackage{tikz}
\usetikzlibrary{positioning,arrows.meta}
\usepackage{url}

% ---------- marcadores de rascunho -----------------------------------
\newcommand{\pendente}[1]{\textcolor{red}{[\textbf{PENDENTE:} #1]}}
\newcommand{\figura}[1]{%
  \fbox{\parbox[c][3.2cm][c]{0.92\columnwidth}{\centering\small\textit{#1}}}}

% ---------- realce de código AEGIS -----------------------------------
\lstdefinelanguage{AEGIS}{
  morekeywords={chronicle,scavenge,mission,supply,law,survivor,cache,stash,
    checkpoint,fallback,watch,patrol,extract,quarantine,catch,outbreak,
    ambush,broadcast,silence,hide,log_entry},
  morekeywords=[2]{ration,vital,status,paper,radio,silent,alive,dead},
  sensitive=true,
  morecomment=[l]{//},
  morecomment=[s]{/*}{*/},
  morestring=[b]",
}
\lstset{
  basicstyle=\ttfamily\scriptsize,
  keywordstyle=\bfseries\color{blue!60!black},
  keywordstyle=[2]\color{teal!70!black},
  commentstyle=\itshape\color{gray},
  stringstyle=\color{red!60!black},
  frame=single,
  breaklines=true,
  columns=fullflexible,
  captionpos=b,
  numbers=left,
  numberstyle=\tiny\color{gray},
  xleftmargin=1.5em,
  framexleftmargin=1.5em,
  literate={á}{{\'a}}1 {é}{{\'e}}1 {í}{{\'i}}1 {ó}{{\'o}}1 {ú}{{\'u}}1
           {ã}{{\~a}}1 {õ}{{\~o}}1 {ç}{{\c{c}}}1 {ê}{{\^e}}1 {â}{{\^a}}1,
}
\renewcommand{\lstlistingname}{Listagem}

\begin{document}

\title{AEGIS: uma Linguagem de Domínio Específico para Automação\\
de Abrigos Residenciais em Microcontroladores}

\author{\IEEEauthorblockN{Yago Duarte de Andrade}
\IEEEauthorblockA{\textit{\pendente{curso}} \\
\textit{\pendente{instituição}}\\
\pendente{cidade}, Brasil \\
\pendente{e-mail}}
}

\maketitle

% =====================================================================
\begin{abstract}
Este artigo apresenta a AEGIS, uma linguagem de domínio específico (DSL)
para descrever as regras de um sistema de monitoramento de abrigos
residenciais em áreas de risco de desastre, e o compilador que a traduz
para Python. O compilador, escrito em JavaScript, é organizado em três
fases: análise léxica por expressões regulares com \textit{maximal munch},
análise sintática descendente recursiva sobre uma gramática LL(1) que
produz uma árvore sintática abstrata (AST) de 26 tipos de nó, e geração de
código com dois alvos, um simulado para PC e outro em MicroPython para o
ESP32. Como estudo de caso, o firmware completo de um abrigo, com sete
sensores, máquina de estados de risco, display OLED, sirene, relé e leitor
biométrico, foi escrito em AEGIS. O compilador passa em 82 testes
automatizados e o firmware em 29 cenários executados no PC e no simulador
Wokwi. \pendente{fechar com uma frase de conclusão}
\end{abstract}

\begin{IEEEkeywords}
compiladores, linguagem de domínio específico, análise sintática LL(1),
sistemas embarcados, ESP32, prevenção de desastres
\end{IEEEkeywords}

% =====================================================================
\section{Introdução}

% Ideia do parágrafo: o problema real (desastres) motiva o produto.
O Brasil tem cerca de 8,27 milhões de pessoas vivendo em áreas de risco de
desastres naturais~\cite{ibge2018}, e mais de 1.600 desastres foram
registrados em 2024~\cite{cemaden2024}. Sistemas comerciais de automação
residencial priorizam o conforto e costumam depender de serviços em
nuvem, o que os torna frágeis justamente quando a energia ou a internet
falham.

% Ideia do parágrafo: por que criar uma linguagem em vez de escrever o
% firmware direto em Python/C.
\pendente{explicar a motivação para uma DSL: as regras de reação a
emergências ficam legíveis como um ``diário de operação'' do abrigo, e o
mesmo programa roda no PC e na placa}

% Ideia do parágrafo: contribuições.
As contribuições deste trabalho são:
\begin{itemize}
  \item a especificação da linguagem AEGIS, com 29 palavras reservadas,
        25 símbolos e uma gramática LL(1) em EBNF;
  \item um compilador de três fases com recuperação de erros e dois alvos
        de geração de código;
  \item o firmware de um abrigo escrito na própria linguagem, validado em
        simulação.
\end{itemize}

O restante do artigo está organizado assim: a Seção~\ref{sec:fundamentos}
revisa os conceitos de linguagens formais usados; a
Seção~\ref{sec:linguagem} apresenta a linguagem; a
Seção~\ref{sec:compilador} descreve o compilador; a
Seção~\ref{sec:firmware} mostra o estudo de caso; a
Seção~\ref{sec:resultados} relata os testes; e a
Seção~\ref{sec:conclusao} conclui.

% =====================================================================
\section{Fundamentação Teórica}
\label{sec:fundamentos}

\subsection{Linguagens regulares e análise léxica}
\pendente{expressões regulares, autômatos finitos, nível 3 da hierarquia
de Chomsky; por que tokens são uma linguagem regular~\cite{hopcroft2006}}

\subsection{Gramáticas livres de contexto e análise sintática}
\pendente{GLC, nível 2 da hierarquia, derivação, árvore de derivação vs.
AST; gramáticas LL(1) e análise descendente recursiva~\cite{aho2007}}

\subsection{Geração de código}
\pendente{tradução dirigida pela sintaxe, visitante sobre a AST,
compilação para outra linguagem de alto nível (\textit{transpiling})}

% =====================================================================
\section{A Linguagem AEGIS}
\label{sec:linguagem}

A AEGIS usa um vocabulário temático de sobrevivência: em vez de
\texttt{function}, \texttt{if} e \texttt{print}, o programador escreve
\texttt{mission}, \texttt{checkpoint} e \texttt{log\_entry}. Os blocos são
delimitados por chaves, como em C, e a quebra de linha ou o ponto e vírgula
encerram um comando. A Tabela~\ref{tab:palavras} resume as palavras
reservadas.

\begin{table}[htbp]
\caption{Palavras reservadas da AEGIS por categoria}
\label{tab:palavras}
\centering
\small
\begin{tabular}{@{}lp{5.2cm}@{}}
\toprule
\textbf{Categoria} & \textbf{Palavras} \\
\midrule
Declaração (8) & \texttt{chronicle scavenge mission supply law survivor cache stash} \\
Tipo (6)       & \texttt{ration vital status paper radio silent} \\
Controle (5)   & \texttt{checkpoint fallback watch patrol extract} \\
Hardware (3)   & \texttt{broadcast silence hide} \\
Erro (3)       & \texttt{quarantine catch outbreak} \\
Evento/E-S (2) & \texttt{ambush log\_entry} \\
Booleanos (2)  & \texttt{alive dead} \\
\bottomrule
\end{tabular}
\end{table}

A Listagem~\ref{lst:exemplo} mostra um programa completo: uma função que
lê um sensor de presença e uma tarefa que aciona um alarme.

\begin{lstlisting}[language=AEGIS,caption={Programa de exemplo em AEGIS},label={lst:exemplo}]
scavenge gpio

law MOTION_PIN = 27
law ALARM_RADIO = 5

mission check_perimeter() -> status {
    supply detected = gpio.read(radio(MOTION_PIN))
    extract detected
}

survivor perimeter_watch {
    watch {
        checkpoint (check_perimeter()) {
            broadcast(radio(ALARM_RADIO))
            hide(10000)
            silence(radio(ALARM_RADIO))
        } fallback {
            log_entry("area limpa, nenhum sinal")
        }
        hide(500)
    }
}
\end{lstlisting}

\pendente{comentar a biblioteca de módulos embutidos (gpio, adc, sonar,
climate, display, wifi, mqtt, finger): não são palavras reservadas, e a
API é a mesma nos dois alvos}

% =====================================================================
\section{O Compilador}
\label{sec:compilador}

O compilador foi escrito em JavaScript sem dependências externas, o que
permite executá-lo tanto pela linha de comando (Node.js) quanto numa
interface web que exibe cada fase. A Fig.~\ref{fig:pipeline} mostra o
encadeamento das fases; a compilação para na primeira fase que acusar
erro.

\begin{figure}[htbp]
\centering
\begin{tikzpicture}[
  node distance=0.35cm,
  fase/.style={draw, rounded corners, fill=blue!8, align=center,
               font=\scriptsize, minimum height=0.9cm, text width=1.35cm},
  dado/.style={align=center, font=\scriptsize\itshape},
  seta/.style={-{Stealth[length=1.6mm]}}]
  \node[dado] (fonte) {fonte\\.aeg};
  \node[fase, right=of fonte] (lex) {Análise\\léxica};
  \node[fase, right=of lex] (sin) {Análise\\sintática};
  \node[fase, right=of sin] (ger) {Geração\\de código};
  \node[dado, right=of ger] (py) {Python};
  \draw[seta] (fonte) -- (lex);
  \draw[seta] (lex) -- node[above, font=\tiny]{tokens} (sin);
  \draw[seta] (sin) -- node[above, font=\tiny]{AST} (ger);
  \draw[seta] (ger) -- (py);
\end{tikzpicture}
\caption{Fases do compilador AEGIS.}
\label{fig:pipeline}
\end{figure}

\subsection{Análise léxica}

O analisador léxico percorre o fonte uma única vez, testando na posição
atual uma lista ordenada de expressões regulares. A primeira que casar
vence, e a ordem implementa a regra do \textit{maximal munch}: símbolos de
dois caracteres antes dos de um (\texttt{->} antes de \texttt{-}),
comentários antes da divisão e literais decimais antes de inteiros. Cada
token registra tipo, categoria, lexema, valor convertido, linha e coluna.

Três erros léxicos são registrados sem interromper a análise: aspas não
fechadas, comentário de bloco não fechado e caractere não reconhecido.
Assim o programador vê todos os erros de uma vez.

\pendente{explicar a quebra de linha dentro de parênteses (como no Python)
e a heurística que evita que um parêntese esquecido engula o resto do
programa}

\figura{Figura: tabela de tokens na interface web\\
(sugestão: print da aba ``Tokens'' do compilador/index.html)}

\subsection{Análise sintática}

O analisador sintático é descendente recursivo, com uma função por regra
da gramática. A gramática é LL(1), com um único ponto que exige olhar dois
tokens à frente: distinguir uma atribuição (\texttt{x = 1}) de uma chamada
solta (\texttt{hide(500)}). A Listagem~\ref{lst:ebnf} mostra a parte da
gramática que define as expressões; a precedência dos operadores está
codificada na ordem das regras, da menor para a maior.

\begin{lstlisting}[caption={Regras de expressão da gramática (EBNF)},label={lst:ebnf},numbers=none]
expressao     = ou ;
ou            = e , { "||" , e } ;
e             = igualdade , { "&&" , igualdade } ;
igualdade     = comparacao , { ("==" | "!=") , comparacao } ;
comparacao    = adicao , { ("<"|">"|"<="|">=") , adicao } ;
adicao        = multiplicacao , { ("+"|"-") , multiplicacao } ;
multiplicacao = unario , { ("*"|"/") , unario } ;
unario        = ("!"|"-") , unario | chamada ;
chamada       = primario , { "(" , [argumentos] , ")"
                           | "." , IDENT | "[" , expressao , "]" } ;
\end{lstlisting}

O resultado é uma AST com 26 tipos de nó, divididos em estrutura do
programa, comandos e expressões. Todo nó guarda linha e coluna, o que
permite apontar erros com precisão.

\pendente{recuperação de erros em modo pânico: o que \texttt{sincronizar()}
descarta e por que sincronizar em qualquer separador}

\pendente{decisões de projeto: palavras chamáveis (\texttt{broadcast},
\texttt{hide}) como identificadores contextuais; encadeamento
\texttt{fallback checkpoint} gerando \texttt{elif}}

\figura{Figura: AST desenhada pela interface web\\
(sugestão: print da aba ``AST'' para o exemplo do checkpoint)}

\subsection{Geração de código}

O gerador é um visitante recursivo sobre a AST. A Tabela~\ref{tab:traducao}
mostra a correspondência das principais construções.

\begin{table}[htbp]
\caption{Tradução de AEGIS para Python}
\label{tab:traducao}
\centering
\small
\begin{tabular}{@{}ll@{}}
\toprule
\textbf{AEGIS} & \textbf{Python} \\
\midrule
\texttt{supply x = 1}                 & \texttt{x = 1} \\
\texttt{mission f(ration a)}          & \texttt{def f(a):} \\
\texttt{checkpoint (c) \{\} fallback} & \texttt{if c: ... else:} \\
\texttt{watch \{\}}                   & \texttt{while True:} \\
\texttt{patrol(n) \{\}}               & \texttt{for \_ in range(n):} \\
\texttt{quarantine \{\} catch (e)}    & \texttt{try: ... except} \\
\texttt{extract x}                    & \texttt{return x} \\
\texttt{\&\& || !}                    & \texttt{and or not} \\
\texttt{alive / dead}                 & \texttt{True / False} \\
\bottomrule
\end{tabular}
\end{table}

Como a linguagem tem construções de hardware que não existem em Python, o
gerador emite um \textit{runtime} embutido no topo do arquivo, e o código
gerado roda sem dependências. Há dois runtimes com a mesma API, escolhidos
pela opção de alvo: \texttt{simulado}, para o PC, com sensores falsos; e
\texttt{esp32}, que usa os módulos \texttt{machine}, \texttt{dht},
\texttt{network} e \texttt{umqtt} do MicroPython. O corpo do programa
traduzido é idêntico nos dois.

\pendente{as outras decisões do gerador: emissão automática de
\texttt{global}, bloco vazio vira \texttt{pass}, mapa de linha gerada
para linha de origem, preservação do tipo \texttt{vital} em \texttt{10.0}}

% =====================================================================
\section{Estudo de Caso: o Firmware do Abrigo}
\label{sec:firmware}

O firmware do abrigo (\texttt{abrigo.aeg}, cerca de 650 linhas) foi
escrito inteiramente em AEGIS e compilado para MicroPython. Um ESP32 lê
presença (PIR), nível da água (HC-SR04), fumaça (MQ-2), qualidade do ar
(MQ-135), temperatura e umidade (DHT22), chuva e umidade do solo. Cada
sensor tem uma tarefa (\texttt{survivor}) que classifica a leitura em
NORMAL, ATENÇÃO ou ALARME; a cada ciclo, o pior risco decide os LEDs, a
sirene, o display OLED, o relé do exaustor e, se houver WiFi, um aviso
por MQTT. O sistema funciona sem internet: a rede é só um canal a mais.

\figura{Figura: diagrama do circuito no Wokwi ou foto da montagem\\
(sugestão: print do firmware/diagram.json aberto no VS Code)}

\pendente{explicar a máquina de estados NORMAL/ATENÇÃO/ALARME e a
prioridade dos riscos (função \texttt{avaliar})}

\pendente{tratamento dos sensores MQ: média móvel, base de ar limpo
adaptativa e histerese; risco de deslizamento combinando chuva e solo}

\pendente{arme e desarme pela digital (AS608 na placa, Touch ID do Mac
na simulação)}

\figura{Figura: telas do display OLED\\
(arquivo pronto em docs/img/telas-oled.png)}

% =====================================================================
\section{Testes e Resultados}
\label{sec:resultados}

O compilador tem uma suíte de 82 testes automatizados, sem dependências
externas, que cobre as três fases e os dois alvos; todos passam. O
firmware é verificado por 29 cenários executados no PC sobre o alvo
simulado, que forçam leituras de sensor e conferem o estado resultante;
todos passam. O circuito também foi executado no simulador
Wokwi~\cite{wokwi}, com as mesmas ligações da placa real, de modo que o
mesmo \texttt{main.py} roda nos dois.

\pendente{tabela com o tamanho de cada módulo do compilador (linhas de
código) e o número de testes por fase}

\pendente{exemplo de mensagem de erro com linha e coluna, mostrando a
recuperação em modo pânico reportando mais de um erro}

\pendente{custo do protótipo (R\$ 235,92 a base; R\$ 408,42 com OLED,
biometria e sensores extras), se o professor quiser a parte de hardware}

% =====================================================================
\section{Conclusão}
\label{sec:conclusao}

\pendente{retomar o objetivo e o que foi entregue}

Como trabalhos futuros, destacam-se: uma fase de análise semântica, já
que os tipos declarados (\texttt{ration}, \texttt{vital}, \texttt{status})
são registrados na AST mas ainda não verificados; alertas pelo Telegram;
um painel web servido pelo próprio ESP32 na rede local; e a montagem do
hardware real.

% =====================================================================
% Referências: confira cada uma antes de entregar.
\begin{thebibliography}{00}
\bibitem{aho2007} A. V. Aho, M. S. Lam, R. Sethi e J. D. Ullman,
\textit{Compilers: Principles, Techniques, and Tools}, 2.\ ed.
Boston: Addison-Wesley, 2007.

\bibitem{hopcroft2006} J. E. Hopcroft, R. Motwani e J. D. Ullman,
\textit{Introduction to Automata Theory, Languages, and Computation},
3.\ ed. Boston: Addison-Wesley, 2006.

\bibitem{ibge2018} IBGE, \textit{População em Áreas de Risco no Brasil}.
Rio de Janeiro: IBGE, 2018. \pendente{conferir dados e link}

\bibitem{cemaden2024} Cemaden, \pendente{relatório de desastres de 2024:
título e link}

\bibitem{wokwi} Wokwi, ``World's most advanced ESP32 simulator''.
[Online]. Disponível: \url{https://wokwi.com}

\bibitem{micropython} MicroPython. [Online]. Disponível:
\url{https://micropython.org} \pendente{citar no texto ou remover}
\end{thebibliography}

\end{document}
